\subsection{RELOJ COMPARADOR (Comparador de cuadrante)}
\label{sec:reloj-comparador}

\subsubsection{Modelo de medición y ley de propagación}

La calibración de un reloj comparador por comparación directa con bloques patrón de
longitud sobre una mesa de planitud de granito \parencite{inacal_pc014, cem_di010}
consiste en leer la desviación que el palpador registra al sustituir la pieza por la
pila de bloques. Sea $\bar{I}$ la indicación promedio del comparador y $L_{BP}$ la
longitud de la pila corregida por certificado. El error de indicación corregido por
efectos térmicos es

\begin{equation}
  E = \bar{I} - L_{BP} + \bar{I}\,\delta\alpha\,\Delta T_i
      + L_{BP}\,\alpha_{BP}\,\delta T,
  \label{eq:reloj-modelo}
\end{equation}

donde $\delta\alpha = \alpha_i - \alpha_{BP}$ es la diferencia entre los coeficientes
de dilatación del instrumento y del bloque patrón, $\Delta T_i = T_i - T_0$ es la
desviación de la temperatura del reloj respecto de la referencia
$T_0 = 20\ ^\circ\text{C}$ y $\delta T = \Delta T_i - \Delta T_{BP}$ es la diferencia
térmica entre ambos cuerpos. Los dos últimos términos son correctivos y de orden
superior frente a $\bar{I} - L_{BP}$.

De la ecuación~\eqref{eq:reloj-modelo} se observa que $E$ es una función diferenciable
de seis variables de entrada,
$E = f(\bar{I}; L_{BP}; \delta\alpha; \Delta T_i; \alpha_{BP}; \delta T)$, por lo que su
varianza se obtiene aplicando la ley de propagación de la
ecuación~\eqref{eq:error-std-values}. Como las variables son estadísticamente
independientes, $\text{Cov}(x_i, x_j) = 0$ para $i \neq j$ y la suma se reduce a

\begin{align}
  \sigma_c^2(E) &=
  \left( \ptij{E}{\bar{I}} \right)^2 \sigma^2(\bar{I}) +
  \left( \ptij{E}{L_{BP}} \right)^2 \sigma^2(L_{BP}) \nonumber \\
  &\quad + \left( \ptij{E}{\delta\alpha} \right)^2 \sigma^2(\delta\alpha) +
  \left( \ptij{E}{\Delta T_i} \right)^2 \sigma^2(\Delta T_i) \label{eq:reloj-gum} \\
  &\quad + \left( \ptij{E}{\alpha_{BP}} \right)^2 \sigma^2(\alpha_{BP}) +
  \left( \ptij{E}{\delta T} \right)^2 \sigma^2(\delta T). \nonumber
\end{align}

Evaluando las derivadas parciales de la función~\eqref{eq:reloj-modelo} respecto de
cada variable de entrada:

\begin{align}
  \ptij{E}{\bar{I}} &= 1 + \delta\alpha \, \Delta T_i, &
  \ptij{E}{L_{BP}} &= -1 + \alpha_{BP} \, \delta T, \nonumber \\
  \ptij{E}{\delta\alpha} &= \bar{I} \, \Delta T_i, &
  \ptij{E}{\Delta T_i} &= \bar{I} \, \delta\alpha, \label{eq:reloj-sens-deriv} \\
  \ptij{E}{\alpha_{BP}} &= L_{BP} \, \delta T, &
  \ptij{E}{\delta T} &= L_{BP} \, \alpha_{BP}. \nonumber
\end{align}

Sustituyendo los coeficientes~\eqref{eq:reloj-sens-deriv} en la
ecuación~\eqref{eq:reloj-gum} se obtiene la varianza combinada del error de
indicación:

\begin{align}
  \sigma_c^2(E) &=
  [1 + \delta\alpha \, \Delta T_i]^2 \sigma^2(\bar{I}) +
  [-1 + \alpha_{BP} \, \delta T]^2 \sigma^2(L_{BP}) \nonumber \\
  &\quad + [\bar{I} \, \Delta T_i]^2 \sigma^2(\delta\alpha) +
  [\bar{I} \, \delta\alpha]^2 \sigma^2(\Delta T_i) \label{eq:reloj-varianza} \\
  &\quad + [L_{BP} \, \delta T]^2 \sigma^2(\alpha_{BP}) +
  [L_{BP} \, \alpha_{BP}]^2 \sigma^2(\delta T). \nonumber
\end{align}

\subsubsection{Componentes de incertidumbre del error de indicación}

Salvo la repetibilidad, que proviene de lecturas repetidas y constituye una evaluación
tipo~A según la ecuación~\eqref{eq:sigma-media-experimental}, las componentes restantes
son evaluaciones tipo~B: derivan de especificaciones del fabricante, del procedimiento
de calibración o de fronteras acotadas conocidas. Para cada una de ellas, acotada en un
intervalo de semiamplitud $a$, se aplica la distribución rectangular de la
ecuación~\eqref{eq:sigma-rectangular}, es decir,

\begin{equation}
  \sigma = \sqrt{\int_{-a}^{a} (x - \bk{x})^2 P(x)\,dx} = \frac{a}{\sqrt{3}}.
  \label{eq:reloj-rectangular}
\end{equation}

\subsubsection*{Repetibilidad de la indicación}

Sea $I_p$ cada una de las $n = 5$ lecturas repetidas en el punto de mayor error
nominal y $\bar{I}$ su media, conforme al procedimiento PC-014
\parencite{inacal_pc014}. La dispersión experimental se estima con la desviación
estándar muestral y la del promedio correspondiente según la
ecuación~\eqref{eq:sigma-media-experimental}:

\begin{equation}
  \sigma(\bar{I}) = \frac{s_R}{\sqrt{n}},
  \qquad
  s_R = \sqrt{\frac{1}{n-1} \sum_{p=1}^{n} (I_p - \bar{I})^2}.
  \label{eq:reloj-sigma-I}
\end{equation}

\subsubsection*{Bloque patrón de referencia}

La longitud de la pila combina la incertidumbre del certificado de calibración con
factor de cobertura $k = 2$ y la deriva dimensional máxima anual estipulada por la
norma ISO~3650 \parencite{iso3650_1998}, esta última modelada mediante la distribución
rectangular:

\begin{equation}
  \sigma^2(L_{BP}) = \left(\frac{U_{cal}}{k}\right)^2 + \frac{U_{der}^2}{3}.
  \label{eq:reloj-sigma-LBP}
\end{equation}

\subsubsection*{Desalineamiento de coseno}

Si el vástago del palpador se inclina un ángulo $\theta \le 1^\circ$ respecto de la
normal del plano de trabajo \parencite{axiospec_dial, cem_di010}, la lectura sobre una
longitud nominal $L$ se acorta en $\delta l_{cos} = L(1 - \cos\theta)$. Tomando esa
desviación como semiamplitud del intervalo,

\begin{equation}
  \sigma(\delta l_{cos}) = \frac{L(1 - \cos\theta)}{2\sqrt{3}}.
  \label{eq:reloj-sigma-cos}
\end{equation}

\subsubsection*{Resolución o división de escala}

Sea $d$ la división de escala de la carátula o pantalla y $m$ el número de partes en
que se subdivide por estimación visual \parencite{jcgm2008gum, axiospec_dial}. El
límite de discriminación es $d/m$, de modo que, con semiamplitud $a = (d/m)/2$,

\begin{equation}
  \sigma(\delta l_{res}) = \frac{d/m}{2\sqrt{3}}.
  \label{eq:reloj-sigma-res}
\end{equation}

\subsubsection*{Planitud de la mesa de soporte}

Dada la tolerancia de planitud $t_{\parallel}$ reportada en el certificado de la mesa
de granito \parencite{cem_di010}, y tomándola como ancho total del intervalo,

\begin{equation}
  \sigma(\delta l_{plan}) = \frac{t_{\parallel}}{2\sqrt{3}}.
  \label{eq:reloj-sigma-plan}
\end{equation}

\subsubsection*{Deformación por fuerza del palpador}

La deformación elástica de un contacto esfera-plano se estima con las relaciones de
Hertz a partir de la fuerza de medida $P$ y del diámetro $D$ de la punta esférica
\parencite{axiospec_dial},

\begin{equation}
  \delta l_{fuerza} = 2\times10^{-5}\, P^{2/3} D^{-1/3}
  \quad\Longrightarrow\quad
  \sigma(\delta l_{fuerza}) = \frac{\delta l_{fuerza}}{2\sqrt{3}}.
  \label{eq:reloj-sigma-fuerza}
\end{equation}

\subsubsection*{Parámetros térmicos}

Sean $\Delta\alpha_i$ y $\Delta\alpha_{BP}$ las semiamplitudes de tolerancia de los
coeficientes de dilatación del instrumento y del bloque. Aplicando
\eqref{eq:reloj-rectangular} a la diferencia $\delta\alpha$,

\begin{equation}
  \sigma^2(\delta\alpha) = \sigma^2(\alpha_i) + \sigma^2(\alpha_{BP})
  = \frac{\Delta\alpha_i^2 + \Delta\alpha_{BP}^2}{3},
  \qquad
  \sigma^2(\alpha_{BP}) = \frac{\Delta\alpha_{BP}^2}{3}.
  \label{eq:reloj-sigma-dalpha}
\end{equation}

Sean $UT_{cert}$ la exactitud del termómetro de contacto, $UT_{der}$ su deriva
máxima, $d_T/m_T$ su límite de discriminación y $\Delta t_A$ la fluctuación observada
durante el ensayo. Combinando las contribuciones independientes con
\eqref{eq:reloj-rectangular},

\begin{equation}
  \sigma^2(\Delta T_i) = \left(\frac{UT_{cert}}{k}\right)^2
  + \left(\frac{d_T/m_T}{\sqrt{3}}\right)^2
  + \frac{UT_{der}^2}{3} + \frac{\Delta t_A^2}{3}.
  \label{eq:reloj-sigma-deltaT}
\end{equation}

De manera análoga, para el gradiente térmico $\delta T$ entre bloque y reloj se define
$\Delta t = \max(\Delta t_1, \Delta t_2)$ como la máxima diferencia térmica entre ambos
cuerpos, de modo que

\begin{equation}
  \sigma^2(\delta T) = \left(\frac{UT_{cert}}{k}\right)^2
  + \frac{\Delta t^2}{3}
  + \left(\frac{d_T/m_T}{\sqrt{3}}\right)^2
  + \frac{UT_{der}^2}{3}.
  \label{eq:reloj-sigma-deltaT-grad}
\end{equation}

\subsubsection*{Varianza combinada de la indicación}

Asumiendo independencia entre las componentes evaluadas, la varianza total de la
indicación es la suma en cuadratura de las
ecuaciones~\eqref{eq:reloj-sigma-I} a~\eqref{eq:reloj-sigma-fuerza}:

\begin{equation}
  \sigma^2(\bar{I}) = \frac{s_R^2}{n}
  + \sigma^2(\delta l_{cos}) + \sigma^2(\delta l_{res})
  + \sigma^2(\delta l_{plan}) + \sigma^2(\delta l_{fuerza}).
  \label{eq:reloj-sigma-I-total}
\end{equation}

\subsubsection{Incertidumbre expandida analítica}

Sustituyendo \eqref{eq:reloj-sigma-I-total} y las
ecuaciones~\eqref{eq:reloj-sigma-LBP} y~\eqref{eq:reloj-sigma-dalpha}
a~\eqref{eq:reloj-sigma-deltaT-grad} en la ley de
propagación~\eqref{eq:reloj-varianza}, la incertidumbre expandida del error de
indicación queda, con factor de cobertura $k$, como

\begin{align}
  U(E) = k \, \Bigg\{ &
  \left[1 + \delta\alpha \, \Delta T_i\right]^2
  \Bigg[ \frac{s_R^2}{n}
    + \left(\frac{L(1-\cos\theta)}{2\sqrt{3}}\right)^2
    + \left(\frac{d/m}{2\sqrt{3}}\right)^2 \nonumber \\
  &\qquad\qquad
    + \left(\frac{t_{\parallel}}{2\sqrt{3}}\right)^2
    + \left(\frac{\delta l_{fuerza}}{2\sqrt{3}}\right)^2 \Bigg] \nonumber \\
  &\quad + \left[-1 + \alpha_{BP} \, \delta T\right]^2
    \left[ \left(\frac{U_{cal}}{k}\right)^2 + \frac{U_{der}^2}{3} \right]
    \nonumber \\
  &\quad + \left[\bar{I} \, \Delta T_i\right]^2
    \left[ \frac{\Delta\alpha_i^2 + \Delta\alpha_{BP}^2}{3} \right]
    + \left[\bar{I} \, \delta\alpha\right]^2 \sigma^2(\Delta T_i) \nonumber \\
  &\quad + \left[L_{BP} \, \delta T\right]^2
    \frac{\Delta\alpha_{BP}^2}{3}
    + \left[L_{BP} \, \alpha_{BP}\right]^2 \sigma^2(\delta T)
  \Bigg\}^{1/2}.
  \label{eq:reloj-U}
\end{align}

Adoptando $k = 2$, correspondiente a un nivel de confianza bilateral aproximado del
$95$\,\% \parencite{jcgm2008gum, inacal_pc014}, la
ecuación~\eqref{eq:reloj-U} queda completamente determinada por las componentes
anteriores y permite evaluar $U(E)$ punto a punto a lo largo del alcance del
instrumento.
